% !PRINTED 327
\weberpage{327}
%<<CONT
nach dem Algorithmus des größten gemeinschaftlichen Teilers dienen. Denn man kann aus den Kongruenzen
\[
\begin{aligned}
D &\equiv D' \pmod{4n}, \\
n &\equiv n' \pmod{D'}, \\
D' &\equiv D'' \pmod{4n'}, \\
n' &\equiv n'' \pmod{D''}, \\
&\cdots\cdots\cdots\cdots
\end{aligned}
\]
die Reihe der Zahlen
\[
D', \quad 2n', \quad D'', \quad 2n'', \quad \dots
\]
so bestimmen, daß jede folgende dem absoluten Werte nach kleiner ist als die vorhergehende, solange keine der Zahlen $n'$, $n''$, $n'''$, $\dots$ gleich $\pm 1$ geworden ist.

Ist $D$ keine Quadratzahl, so läßt sich immer eine Zahl $\beta$ so bestimmen, daß
\weq{21}{(D, \beta) = -1.}

Ist nämlich $\varDelta$ der Stamm von $D$ und ist $\beta$ relativ prim zu $D$, so ist (nach 9)
\[
(D, \beta) = (\varDelta, \beta).
\]

Ist $\varDelta = \delta \varDelta'$ und $\delta$ eine Primdiskriminante, also relativ prim zu $\varDelta'$, so kann man $\beta_0$ so annehmen, daß
\[
(\delta, \beta_0) = -1
\]
wird [nach (10) und (11)] und man kann $\beta$ aus den beiden Kongruenzen
\[
\begin{aligned}
\beta &\equiv \beta_0 \pmod{\delta}, \\
&\equiv 1 \pmod{\varDelta'}
\end{aligned}
\]
bestimmen. Dann erhält man nach (12)
\[
(D, \beta) = (\delta, \beta_0) = -1.
\]

Läßt man $s$ ein volles Restsystem nach dem Modul $D$ durchlaufen und setzt
\[
S = \overset{s}{\Sigma}(D, s), % Source scan confirms s as the upper index on Sigma.
\]
so folgt durch Multiplikation mit $(D, \beta)$
\[
-S = \overset{s}{\Sigma}(D, \beta s),
\]
und da $\beta s$ zugleich mit $s$ ein Restsystem durchläuft, so folgt
\weq{22}{S = 0.}

Es folgt hieraus:
%CONT>>
